\begin{center}
\SourceLeaf{030b}\pdfbookmark[2]{12.}{article-012}12.
\end{center}
Valor approximatus ipsius $E$, a quo calculus incipi possit, plerumque satis obvius erit, praesertim ubi problema pro pluribus valoribus ipsius $M$ solvendum est, e quibus quidam iam absoluti sunt. Deficientibus omnibus aliis subsidiis id saltem constat, quod $E$ inter limites $M$ et $M \pm e$ iacere debet (excentricitate $e$ in secundis expressa, signoque superiori in quadrante primo et secundo, inferiori in tertio et quarto accepto); quocirca pro valore initiali ipsius $E$ vel $M$ vel valor secundum aestimationem qualemcunque auctus seu deminutus adoptari poterit. Vix opus est monere, calculum primum, quoties a valore parum accurato inchoetur, anxia praecisione haud indigere, tabulasque minores quales cel. \textsc{Lalande} curavit, abunde sufficere. Praeterea, ut calculi commoditati consulatur, tales semper valores pro $\epsilon$ eligentur, quorum sinus e tabulis ipsis absque interpolatione excerpere licet; puta in minutis seu secundorum denariis completis, prout tabulae per singula minuta seu per singulos secundorum denarios progredientes adhibentur. Ceterum modificationes, quas haec praecepta patiuntur, si anguli secundum divisionem novam decimalem exprimantur, quisque sponte evolvere poterit.
